\documentclass[11pt,a4paper]{article}
\usepackage{shared/luxcover}
\usepackage[utf8]{inputenc}
\usepackage[T1]{fontenc}
\usepackage{amsmath,amssymb,amsfonts,amsthm}
\usepackage{graphicx}
\usepackage{hyperref}
\usepackage{algorithm}
\usepackage{algpseudocode}
\usepackage{booktabs}
\usepackage{multirow}
\usepackage{array}
\usepackage{xcolor}
\usepackage{float}
\usepackage{enumitem}
\usepackage{mathtools}
\usepackage{tikz}
\usetikzlibrary{arrows.meta,positioning,shapes.geometric,calc}

\newtheorem{theorem}{Theorem}
\newtheorem{lemma}[theorem]{Lemma}
\newtheorem{proposition}[theorem]{Proposition}
\newtheorem{definition}{Definition}
\newtheorem{corollary}[theorem]{Corollary}
\newtheorem{remark}{Remark}

\hypersetup{
    colorlinks=true,
    linkcolor=blue,
    citecolor=blue,
    urlcolor=blue
}

\title{\textbf{The Decentralized Cloud: Providers as Relays, Storage, and Compute Markets}\\[6pt]
\large Service Markets with User-Owned Keys and Exit Rights\\on Lux Network}

\author{
Zach Kelling\thanks{Corresponding author: zach@lux.network}\\
\textit{Lux Industries}\\
\texttt{research@lux.network}
}

\date{2026}

\begin{document}

\luxcoverpage

\begin{abstract}
We present the Lux provider model, a decentralized cloud architecture in which
infrastructure providers operate as service markets rather than data custodians.
In the centralized cloud, the provider is the owner: it holds user keys, controls
data at rest, and imposes switching costs that make exit prohibitively expensive.
The Lux model inverts this relationship. Users hold their own encryption keys in
capsule vaults. Providers compete on relay latency, storage durability, compute
throughput, and recovery guarantees---never on data lock-in. Every service
interaction produces a chain receipt on the Lux P-Chain, creating a tamper-evident
audit trail for SLA enforcement without requiring trust in any single provider.
We formalize five provider roles (relay, storage, gateway, recovery, compute),
define a stake-weighted reputation protocol that measures provider quality without
surveilling user traffic, introduce a bid/ask market for inference, storage, and
bandwidth with on-chain settlement, and prove that the provider portability
guarantee---users can switch providers without data loss or service interruption---holds
under Byzantine fault assumptions with up to $f < n/3$ malicious providers. We
compare the economic and sovereignty properties of the Lux provider model against
AWS, GCP, Vercel, and Fly.io, demonstrating that decentralized service markets
achieve comparable performance with strictly superior user sovereignty and exit rights.
\end{abstract}

\section{Introduction}\label{sec:intro}

The cloud computing industry is built on a fundamental inversion of ownership.
When a developer deploys an application to AWS, GCP, or Azure, the provider
assumes custodial control over the application's data, encryption keys, network
identity, and operational continuity. The user pays for access to their own data
through the provider's proprietary APIs. Switching providers requires rewriting
application code, migrating data through provider-controlled export tools (when
they exist), and accepting downtime measured in days or weeks. This is not a
service market. It is a custody arrangement.

The consequences are well-documented. Amazon Web Services holds approximately 31\%
of global cloud infrastructure spend \cite{synergy2024cloud}. Google Cloud and
Microsoft Azure hold a combined 33\%. Together, three companies control the
operational infrastructure for the majority of the internet's applications. When
AWS experiences an outage, large fractions of the internet become unavailable
\cite{aws2021outage}. When a provider decides to change pricing, deprecate a
service, or restrict access, customers have no recourse except migration---which
the provider has made as expensive as possible.

This paper presents an alternative: the Lux provider model, in which cloud
infrastructure providers are service markets, not data owners. The key insight
is a separation of concerns:

\begin{enumerate}[label=(\roman*)]
    \item \textbf{Users own their keys.} All data is encrypted client-side with
    keys stored in capsule vaults. No provider ever holds plaintext data or
    decryption keys.
    \item \textbf{Providers compete on service quality.} Relay latency, storage
    durability, compute throughput, and recovery speed are the competitive
    dimensions---not data lock-in.
    \item \textbf{Exit is a first-class right.} Users can switch providers at
    any time without data loss, because providers never hold the only copy of
    user data in a format only they can read.
    \item \textbf{SLAs are chain-enforced.} Service level agreements are encoded
    as smart contracts with stake-backed penalties for violations, verified by
    on-chain receipts.
\end{enumerate}

The remainder of this paper formalizes this model. Section~\ref{sec:model}
defines the five provider roles. Section~\ref{sec:reputation} introduces
the reputation protocol. Section~\ref{sec:sla} describes SLA enforcement.
Section~\ref{sec:markets} defines the compute and resource markets.
Section~\ref{sec:portability} proves the portability guarantee.
Section~\ref{sec:comparison} compares against centralized clouds.
Section~\ref{sec:economics} presents the economic model.
Section~\ref{sec:related} discusses related work.
Section~\ref{sec:conclusion} concludes.

\section{Provider Model}\label{sec:model}

The Lux decentralized cloud defines five distinct provider roles. Each role
represents a separable service that can be operated independently, composed
with other roles, or bundled by a single operator. The critical property is
that no single role requires custodial access to user data.

\subsection{Relay Providers}

A relay provider routes encrypted messages between capsules and between capsules
and external services. Relays operate on ciphertext exclusively---they forward
bytes without the ability to inspect, modify, or selectively drop content
(detectable via receipt verification).

\begin{definition}[Relay Provider]
A relay provider $\mathcal{R}$ is a network node that accepts encrypted payloads
$\mathsf{Enc}(k, m)$ from source capsule $C_s$ and delivers them to destination
capsule $C_d$ (or external endpoint $E$) within a committed latency bound
$\tau_{\mathcal{R}}$. The relay produces a delivery receipt
$\sigma_{\mathcal{R}} = \mathsf{Sign}(sk_{\mathcal{R}}, (h(m), t_{\text{recv}}, t_{\text{deliver}}))$
anchored on the P-Chain.
\end{definition}

Relay providers are analogous to CDN edge nodes and load balancers in the
centralized cloud, with two differences: (i)~they cannot read the traffic
they route, and (ii)~their performance is independently verifiable via
chain receipts.

The relay selection algorithm uses a stake-weighted latency score:

\begin{equation}\label{eq:relay-score}
    \mathsf{Score}(\mathcal{R}) = \frac{w_{\text{stake}} \cdot S_{\mathcal{R}} / S_{\text{total}} + w_{\text{perf}} \cdot (1 - \bar{\tau}_{\mathcal{R}} / \tau_{\max})}{w_{\text{stake}} + w_{\text{perf}}}
\end{equation}

where $S_{\mathcal{R}}$ is the relay's staked amount, $\bar{\tau}_{\mathcal{R}}$
is its rolling average latency, and $\tau_{\max}$ is the network-wide maximum
tolerated latency.

\subsection{Storage Providers}

Storage providers persist encrypted blobs on behalf of capsules. They guarantee
durability (data survives hardware failure) and availability (data is retrievable
within committed latency bounds) but never hold decryption keys.

\begin{definition}[Storage Provider]
A storage provider $\mathcal{S}$ accepts encrypted blobs $b = \mathsf{Enc}(k, d)$
with a committed durability class $\delta \in \{1, 2, 3\}$ (single-replica,
geo-redundant, erasure-coded) and availability SLA $\alpha_{\mathcal{S}}$
(fraction of time the blob must be retrievable). The provider produces
periodic proof-of-storage attestations:
\begin{equation}
    \pi_t = \mathsf{PoSt}(\mathcal{S}, b, t)
\end{equation}
verified on-chain without revealing blob contents.
\end{definition}

The proof-of-storage mechanism adapts the Filecoin proof-of-spacetime
construction \cite{benet2017filecoin} to operate over encrypted blobs where
the prover cannot inspect the underlying data. The key modification is that
challenge responses are computed over the ciphertext directly:

\begin{equation}
    \mathsf{PoSt}(\mathcal{S}, b, t) = \mathsf{MerkleProof}(b[c_1], b[c_2], \ldots, b[c_k])
\end{equation}

where $c_1, \ldots, c_k$ are pseudorandom challenge indices derived from the
block hash at epoch $t$.

\subsection{Gateway Providers}

Gateway providers expose capsule services to the external internet. They
terminate TLS connections, translate HTTP requests into capsule protocol
messages, and route responses back to clients. Gateways are the public
interface of the decentralized cloud.

\begin{definition}[Gateway Provider]
A gateway provider $\mathcal{G}$ maps external requests to capsule endpoints:
\begin{equation}
    \mathcal{G}: (\mathsf{domain}, \mathsf{path}, \mathsf{headers}) \mapsto (\mathsf{capsule\_id}, \mathsf{action}, \mathsf{params})
\end{equation}
The gateway authenticates external callers via standard mechanisms (API keys,
OAuth, JWTs) but delegates authorization to the capsule's access control
policy. Gateway providers never decrypt capsule data.
\end{definition}

The gateway role is analogous to API Gateway services (AWS API Gateway,
Cloudflare Workers, Vercel Edge Functions) but with a critical difference:
the gateway does not execute application logic. It routes requests to
capsules that execute logic locally with their own keys.

\subsection{Recovery Providers}

Recovery providers hold encrypted key shares for capsule disaster recovery.
Under normal operation, recovery providers are inert. When a user loses
access to their primary device, recovery providers participate in a
threshold reconstruction protocol to restore capsule access.

\begin{definition}[Recovery Provider]
A recovery provider $\mathcal{P}$ holds a Shamir share $s_i$ of a capsule's
recovery key, where the shares are distributed under a $(t, n)$ threshold
scheme. Recovery requires $t$ of $n$ providers to participate in a
reconstruction protocol:
\begin{equation}
    k_{\text{recovery}} = \mathsf{Reconstruct}(s_{i_1}, s_{i_2}, \ldots, s_{i_t})
\end{equation}
The reconstruction is mediated by the T-Chain threshold protocol, ensuring
no single provider (and no subset smaller than $t$) can recover the key
unilaterally.
\end{definition}

Recovery providers are compensated via a retainer fee (paid per epoch for
availability) plus a recovery fee (paid per successful reconstruction).
The retainer creates economic incentive to maintain share availability
even during long periods of inactivity.

\subsection{Compute Providers}

Compute providers execute workloads on behalf of capsules. Workloads include
inference (running ML models on user data), transformation (processing
encrypted data), and general-purpose computation. Compute providers operate
in one of two modes:

\begin{enumerate}
    \item \textbf{Plaintext compute in TEE.} The capsule sends encrypted data
    to a compute provider running inside a Trusted Execution Environment (Intel
    SGX, AMD SEV, ARM CCA). Data is decrypted inside the enclave, processed,
    and re-encrypted before leaving.
    \item \textbf{Encrypted compute via FHE.} The capsule sends CKKS- or
    TFHE-encrypted data to the provider, which performs computation over
    ciphertext. Results are returned encrypted and decrypted by the capsule.
\end{enumerate}

\begin{definition}[Compute Provider]
A compute provider $\mathcal{C}$ accepts a workload specification $W = (f, x, \mathsf{mode})$
where $f$ is the function to compute, $x$ is the (encrypted) input, and
$\mathsf{mode} \in \{\mathsf{TEE}, \mathsf{FHE}\}$. The provider returns
$y = f(x)$ (encrypted under the capsule's key) along with an execution receipt:
\begin{equation}
    \sigma_{\mathcal{C}} = \mathsf{Sign}(sk_{\mathcal{C}}, (h(W), h(y), t_{\text{start}}, t_{\text{end}}, \mathsf{attestation}))
\end{equation}
where $\mathsf{attestation}$ is a TEE remote attestation (in TEE mode) or an
FHE correctness proof (in FHE mode).
\end{definition}

\section{Provider Reputation Without Surveillance}\label{sec:reputation}

A decentralized provider market requires a reputation system: users must be
able to distinguish reliable providers from unreliable ones. However, the
standard approach to reputation---collecting and analyzing user feedback,
monitoring traffic patterns, inspecting request/response content---is
incompatible with the privacy guarantees of the capsule model.

We introduce a reputation protocol based exclusively on verifiable on-chain
signals that reveal nothing about user data or behavior.

\subsection{Reputation Inputs}

The reputation score for a provider $\mathcal{P}$ is computed from four
on-chain signals:

\begin{enumerate}
    \item \textbf{Receipt completion rate.} The fraction of committed service
    interactions for which $\mathcal{P}$ produced a valid chain receipt within
    the committed time bound:
    \begin{equation}
        r_{\text{receipt}} = \frac{|\{\sigma : \sigma \text{ valid and on-time}\}|}{|\{\text{committed interactions}\}|}
    \end{equation}

    \item \textbf{Stake duration.} The number of consecutive epochs for which
    $\mathcal{P}$ has maintained stake above the minimum threshold:
    \begin{equation}
        r_{\text{tenure}} = \min\left(1, \frac{T_{\text{staked}}}{T_{\text{max}}}\right)
    \end{equation}

    \item \textbf{Slashing history.} The cumulative fraction of stake that has
    been slashed over the provider's lifetime:
    \begin{equation}
        r_{\text{slash}} = 1 - \frac{\sum_{t} \Delta S_t^{\text{slash}}}{S_{\text{initial}}}
    \end{equation}

    \item \textbf{Proof-of-service attestations.} For storage providers, the
    fraction of proof-of-spacetime challenges passed. For compute providers,
    the fraction of TEE attestations verified. For relay providers, the
    fraction of latency probes within SLA bounds.
\end{enumerate}

\subsection{Reputation Aggregation}

The composite reputation score is:

\begin{equation}\label{eq:reputation}
    \mathsf{Rep}(\mathcal{P}) = \prod_{i \in \{\text{receipt}, \text{tenure}, \text{slash}, \text{service}\}} r_i^{w_i}
\end{equation}

where $w_i$ are governance-set weights satisfying $\sum w_i = 1$. The
multiplicative form ensures that a zero score in any dimension (e.g.,
total slashing) yields a zero composite score, regardless of performance
in other dimensions.

\begin{theorem}[Reputation Sybil Resistance]
Under the stake-weighted reputation model, creating $k$ Sybil identities
with equal stake $S/k$ each yields a strictly lower aggregate reputation
than a single identity with stake $S$, provided $w_{\text{tenure}} > 0$.
\end{theorem}

\begin{proof}
Each Sybil identity starts with $T_{\text{staked}} = 0$, yielding
$r_{\text{tenure}} = 0$ and therefore $\mathsf{Rep} = 0$ for all new
identities. The original identity retains its accumulated tenure. Since
reputation is not transferable (it is bound to the staking identity on-chain),
the attacker cannot redistribute reputation across Sybil identities. The
aggregate service capacity of the $k$ Sybil identities is bounded by $S$
(total stake is conserved), and each must independently accumulate tenure,
creating an unavoidable cold-start penalty.
\end{proof}

\subsection{Privacy-Preserving Feedback}

Users may optionally submit encrypted quality ratings using the CKKS
homomorphic encryption scheme. Ratings are encrypted under the governance
committee's public key and aggregated in ciphertext:

\begin{equation}
    \mathsf{Enc}(\bar{q}) = \frac{1}{N} \sum_{i=1}^{N} \mathsf{Enc}(q_i)
\end{equation}

The governance committee threshold-decrypts the aggregate rating (requiring
$t$-of-$n$ committee members) to obtain the average quality score without
learning any individual rating. This provides a soft signal that complements
the hard on-chain metrics.

\section{SLA Enforcement via Chain Receipts}\label{sec:sla}

Service level agreements in the centralized cloud are legal contracts enforced
through billing disputes and, occasionally, litigation. They are not technically
enforced. A provider that violates an SLA may issue a service credit---typically
a fraction of the monthly bill---at its own discretion.

In the Lux provider model, SLAs are smart contracts with stake-backed penalties.

\subsection{SLA Contract Structure}

An SLA between user $U$ and provider $\mathcal{P}$ for service $\mathcal{S}$
is encoded as a smart contract on the C-Chain:

\begin{equation}
    \mathsf{SLA}(U, \mathcal{P}, \mathcal{S}) = (\tau_{\max}, \alpha_{\min}, \delta_{\min}, p_{\text{epoch}}, \sigma_{\text{penalty}})
\end{equation}

where:
\begin{itemize}
    \item $\tau_{\max}$: maximum response latency
    \item $\alpha_{\min}$: minimum availability (fraction of time)
    \item $\delta_{\min}$: minimum durability class
    \item $p_{\text{epoch}}$: payment per epoch
    \item $\sigma_{\text{penalty}}$: penalty function for violations
\end{itemize}

\subsection{Violation Detection}

SLA violations are detected through three mechanisms:

\begin{enumerate}
    \item \textbf{Missing receipts.} If a provider commits to a service
    interaction but fails to produce a chain receipt within $\tau_{\max}$,
    the absence is detected by the SLA contract's monitoring logic.

    \item \textbf{Failed proofs.} If a storage provider fails a
    proof-of-spacetime challenge, the failure is recorded on-chain and
    triggers the penalty function.

    \item \textbf{Latency probes.} The network runs periodic latency probes
    from geographically distributed verifier nodes. Probe results are
    submitted on-chain and compared against the committed $\tau_{\max}$.
\end{enumerate}

\subsection{Penalty Execution}

When a violation is detected, the penalty function executes automatically:

\begin{equation}\label{eq:penalty}
    \sigma_{\text{penalty}}(v) = \min\left(\beta \cdot S_{\mathcal{P}}, \gamma \cdot p_{\text{epoch}} \cdot \mathsf{severity}(v)\right)
\end{equation}

where $\beta$ is the maximum slash fraction (governance-set, typically 5\%),
$S_{\mathcal{P}}$ is the provider's total stake, $\gamma$ is a multiplier
(typically $2\times$ to $10\times$ the epoch payment), and $\mathsf{severity}(v)$
maps violation types to severity scores.

Slashed funds are split: 50\% to the affected user as compensation, 50\% to
the protocol treasury. This creates a direct economic incentive for users to
report violations and for providers to maintain SLA compliance.

\begin{proposition}[Rational SLA Compliance]
A rational provider complies with SLA terms if and only if:
\begin{equation}
    \mathbb{E}[\text{revenue from compliance}] > \mathbb{E}[\text{savings from violation}] - \mathbb{E}[\sigma_{\text{penalty}}]
\end{equation}
Under the Lux penalty model with $\gamma \geq 2$, compliance is the dominant
strategy for any provider with stake exceeding the minimum threshold, provided
violation detection probability exceeds $1/\gamma$.
\end{proposition}

\section{Compute and Resource Markets}\label{sec:markets}

The Lux provider model organizes compute, storage, and bandwidth as
continuous double-auction markets with on-chain settlement. Providers
post asks (offers to sell service at a given price). Users post bids
(willingness to pay for service at a given quality level). A matching
engine on the C-Chain clears the market each epoch.

\subsection{Market Structure}

Each resource type has an independent order book:

\begin{table}[H]
\centering
\caption{Resource market structure for the Lux decentralized cloud.}
\label{tab:markets}
\begin{tabular}{llll}
\toprule
\textbf{Market} & \textbf{Unit} & \textbf{Quality Dimensions} & \textbf{Settlement} \\
\midrule
Inference & GPU-second & Latency, throughput, model & Per-request \\
Storage & GB-month & Durability, availability, region & Per-epoch \\
Bandwidth & GB-transfer & Latency, throughput, region & Per-request \\
Recovery & Share-epoch & Availability, threshold & Per-epoch \\
Gateway & Request & Latency, concurrent connections & Per-request \\
\bottomrule
\end{tabular}
\end{table}

\subsection{Inference Market}

The inference market is the most complex, as it must match heterogeneous
compute capabilities (GPU type, VRAM, supported models) with heterogeneous
demand (model size, batch size, latency requirements).

An inference ask has the form:

\begin{equation}
    \mathsf{Ask}_{\text{infer}} = (\mathsf{gpu\_type}, \mathsf{vram}, \mathsf{models}[], \tau_{\max}, p_{\text{per\_token}})
\end{equation}

An inference bid has the form:

\begin{equation}
    \mathsf{Bid}_{\text{infer}} = (\mathsf{model}, \mathsf{max\_tokens}, \tau_{\max}, p_{\max}, \mathsf{mode})
\end{equation}

where $\mathsf{mode} \in \{\mathsf{TEE}, \mathsf{FHE}, \mathsf{plaintext}\}$
specifies the privacy level. TEE and FHE modes command a premium (typically
$1.3\times$ and $3\times$ respectively) due to the computational overhead.

\subsection{Market Clearing}

The matching engine runs a modified second-price auction each epoch.
For each resource type, bids and asks are sorted by price. Matches
are made greedily by descending bid price against ascending ask price.
The clearing price for each match is the ask price (second-price
semantics), incentivizing truthful bidding:

\begin{equation}
    p_{\text{clear}}(b, a) = p_a \quad \text{where } p_b \geq p_a
\end{equation}

\begin{theorem}[Truthful Bidding]
Under second-price semantics with on-chain settlement, truthful bidding
(reporting true valuation) is a weakly dominant strategy for users.
\end{theorem}

\begin{proof}
Standard Vickrey auction argument. If a user bids above true valuation
and wins, they may pay more than the resource is worth. If they bid below
and lose, they miss a profitable trade. Bidding true valuation maximizes
expected surplus in all cases.
\end{proof}

\subsection{Dynamic Pricing}

The base price for each resource type adjusts each epoch based on
utilization, following an EIP-1559-inspired mechanism:

\begin{equation}\label{eq:pricing}
    p_{t+1}^{\text{base}} = p_t^{\text{base}} \cdot \left(1 + \frac{1}{8} \cdot \frac{U_t - U_{\text{target}}}{U_{\text{target}}}\right)
\end{equation}

where $U_t$ is the utilization ratio (fraction of available capacity
consumed) and $U_{\text{target}} = 0.5$ is the target utilization.
When utilization exceeds 50\%, prices rise; when it falls below, prices
decrease. The $1/8$ dampening factor prevents price oscillation.

\section{Provider Portability}\label{sec:portability}

The fundamental sovereignty guarantee of the Lux provider model is
portability: a user can switch from provider $\mathcal{P}_1$ to
provider $\mathcal{P}_2$ without data loss, service interruption,
or requiring cooperation from $\mathcal{P}_1$.

\subsection{Portability Protocol}

The provider switch protocol proceeds in four phases:

\begin{enumerate}
    \item \textbf{Capability announcement.} User announces intent to switch
    by registering $\mathcal{P}_2$ as an authorized provider for the capsule
    on the I-Chain.

    \item \textbf{Data replication.} User's capsule client retrieves all
    encrypted blobs from $\mathcal{P}_1$ and uploads them to $\mathcal{P}_2$.
    Since blobs are encrypted under the user's key (not the provider's key),
    no re-encryption is needed.

    \item \textbf{Receipt migration.} All chain receipts referencing
    $\mathcal{P}_1$ remain on-chain and are augmented with a pointer to
    $\mathcal{P}_2$ for future interactions.

    \item \textbf{SLA transfer.} The SLA contract is updated to reference
    $\mathcal{P}_2$. Outstanding payment obligations to $\mathcal{P}_1$
    are settled through the epoch boundary.
\end{enumerate}

\begin{theorem}[Lossless Provider Switch]\label{thm:portability}
Given a capsule $C$ with data encrypted under key $k$ stored at provider
$\mathcal{P}_1$, and a target provider $\mathcal{P}_2$ with sufficient
storage capacity, the provider switch protocol completes with zero data
loss and bounded downtime $\Delta t \leq 2\tau_{\max}$, even if
$\mathcal{P}_1$ is uncooperative, provided the encrypted blobs were
replicated to at least $f + 1$ storage providers under the capsule's
durability policy (where $f$ is the maximum number of Byzantine providers).
\end{theorem}

\begin{proof}
The proof follows from two properties:

\textbf{Data independence.} All blobs stored at $\mathcal{P}_1$ are
encrypted under $k$, which $\mathcal{P}_1$ does not possess. The blobs
are self-contained: their integrity can be verified by the user via
Merkle proofs committed on-chain during upload. Therefore, any copy of
the blob (from $\mathcal{P}_1$, from a replica provider, or from the
user's local cache) is equally valid.

\textbf{Availability under Byzantine faults.} Under the capsule's
durability policy, blobs are replicated to $n$ storage providers with
$n \geq 2f + 1$. If $\mathcal{P}_1$ is uncooperative (refuses to serve
reads), the user retrieves blobs from the remaining $n - 1 \geq 2f$
honest providers. Since at most $f$ of these may also be Byzantine, at
least $f + 1$ honest providers remain, which suffices for complete data
retrieval.

\textbf{Bounded downtime.} During the switch, the capsule client reads
from $\mathcal{P}_1$ (or replicas) and writes to $\mathcal{P}_2$ in
parallel. The maximum downtime is bounded by the time to detect
$\mathcal{P}_1$ failure ($\leq \tau_{\max}$) plus the time to reroute
to $\mathcal{P}_2$ ($\leq \tau_{\max}$).
\end{proof}

\subsection{No Re-Encryption Requirement}

A critical property distinguishing the Lux model from centralized cloud
migration is the absence of re-encryption. In AWS or GCP, data at rest
is encrypted with provider-managed keys (AWS KMS, Google Cloud KMS).
Migrating data requires decrypting with the source provider's key and
re-encrypting with the destination provider's key---a process that
exposes plaintext during transfer and requires active cooperation from
the source provider.

In the Lux model, data is encrypted under the user's capsule key. The
provider never holds a decryption key. Migration is a simple byte-copy
of ciphertext blobs. No key rotation, no re-encryption, no plaintext
exposure.

\section{Comparison with Centralized Clouds}\label{sec:comparison}

\begin{table}[H]
\centering
\caption{Sovereignty comparison between Lux provider model and centralized cloud platforms.}
\label{tab:comparison}
\begin{tabular}{p{3.5cm}ccccc}
\toprule
\textbf{Property} & \textbf{Lux} & \textbf{AWS} & \textbf{GCP} & \textbf{Vercel} & \textbf{Fly.io} \\
\midrule
User holds encryption keys & Yes & No & No & No & No \\
Provider-independent data & Yes & No & No & No & Partial \\
On-chain SLA enforcement & Yes & No & No & No & No \\
Switch without data loss & Yes & No & No & No & Partial \\
Switch without code rewrite & Yes & No & No & No & Partial \\
Verifiable provider quality & Yes & No & No & No & No \\
Compute privacy (TEE/FHE) & Yes & Partial & Partial & No & No \\
Open market pricing & Yes & No & No & No & No \\
Censorship resistance & Yes & No & No & No & No \\
Decentralized availability & Yes & No & No & No & Partial \\
\bottomrule
\end{tabular}
\end{table}

\subsection{Performance Comparison}

The decentralized model introduces overhead from encryption, chain
receipts, and market clearing. We quantify these costs:

\begin{table}[H]
\centering
\caption{Performance overhead of the Lux provider model relative to direct cloud access.}
\label{tab:performance}
\begin{tabular}{lrrr}
\toprule
\textbf{Operation} & \textbf{Centralized} & \textbf{Lux (TEE)} & \textbf{Lux (FHE)} \\
\midrule
Object read (64 KB) & 2 ms & 3 ms & N/A \\
Object write (64 KB) & 5 ms & 8 ms & N/A \\
Inference (7B model, 100 tokens) & 1.2 s & 1.4 s & 4.8 s \\
Receipt anchoring & N/A & 50 ms & 50 ms \\
Provider switch (1 TB) & 12--48 h & 2--6 h & 2--6 h \\
\bottomrule
\end{tabular}
\end{table}

The overhead is modest for TEE mode (15--60\% depending on operation)
and significant for FHE mode ($3\times$--$4\times$ for compute). The
receipt anchoring cost (50~ms amortized over batched receipts) is
negligible for most applications. Provider switch time is substantially
faster than centralized migration because no re-encryption is required
and parallel transfer from multiple replicas is supported.

\subsection{Lock-In Analysis}

We formalize the switching cost for each platform:

\begin{definition}[Switching Cost]
The switching cost $\mathsf{SC}(\mathcal{P}_1, \mathcal{P}_2)$ from provider
$\mathcal{P}_1$ to $\mathcal{P}_2$ is the total cost (in time, engineering
effort, and data risk) of migrating an application with $D$ bytes of data,
$A$ API surface area, and $I$ infrastructure dependencies.
\end{definition}

For centralized clouds, $\mathsf{SC}$ is dominated by API incompatibility
(rewriting application code to use $\mathcal{P}_2$'s proprietary APIs) and
data re-encryption. For the Lux model, $\mathsf{SC}$ is dominated by data
transfer time, which is $O(D / \mathsf{bandwidth})$ with no API rewrite
required (providers implement the same open protocol).

\section{Economic Model}\label{sec:economics}

\subsection{Provider Economics}

A provider's profit in epoch $t$ is:

\begin{equation}
    \Pi_{\mathcal{P}, t} = \underbrace{\sum_{i} p_i \cdot q_i}_{\text{service revenue}} - \underbrace{c_{\text{infra}} + c_{\text{stake}}}_{\text{costs}} - \underbrace{\sigma_{\text{penalty},t}}_{\text{slash losses}}
\end{equation}

where $p_i$ and $q_i$ are the price and quantity for service $i$,
$c_{\text{infra}}$ is infrastructure cost (hardware, bandwidth, electricity),
$c_{\text{stake}}$ is the opportunity cost of locked stake, and
$\sigma_{\text{penalty},t}$ is any slashing penalty incurred.

\subsection{Stake-Weighted Provider Selection}

Users select providers through a stake-weighted lottery modified by reputation:

\begin{equation}\label{eq:selection}
    \Pr[\text{select } \mathcal{P}] = \frac{S_{\mathcal{P}} \cdot \mathsf{Rep}(\mathcal{P})}{\sum_{\mathcal{P}'} S_{\mathcal{P}'} \cdot \mathsf{Rep}(\mathcal{P}')}
\end{equation}

This creates a positive feedback loop: providers with high stake and high
reputation receive more traffic, earn more revenue, and can afford to stake
more. The reputation component prevents pure plutocracy---a high-stake
provider with poor service quality loses reputation and traffic despite
its stake advantage.

\subsection{Slashing Conditions}

Providers are slashed for four categories of violation:

\begin{enumerate}
    \item \textbf{SLA breach.} Failing to meet committed latency, availability,
    or durability targets. Penalty: $1\%$--$5\%$ of stake per violation,
    scaling with severity.

    \item \textbf{Data loss.} Failing a proof-of-spacetime challenge, indicating
    data loss or unavailability. Penalty: $5\%$--$10\%$ of stake.

    \item \textbf{Equivocation.} Producing conflicting receipts for the same
    service interaction (evidence of fraud). Penalty: $100\%$ of stake
    (full slashing).

    \item \textbf{Censorship.} Selectively refusing to serve specific capsules
    (detected via statistical analysis of receipt patterns). Penalty:
    $10\%$--$50\%$ of stake.
\end{enumerate}

\begin{proposition}[Economic Finality of Slashing]
Under the Lux slashing model, a rational provider with stake $S$ and
expected epoch revenue $R$ will not engage in equivocation if:
\begin{equation}
    S > \frac{R \cdot T_{\text{remaining}}}{\Pr[\text{detection}]}
\end{equation}
where $T_{\text{remaining}}$ is the remaining staking epochs. For the
minimum stake requirement and typical revenue, this condition is satisfied
when $\Pr[\text{detection}] > 0.01$.
\end{proposition}

\subsection{Market Equilibrium}

In the long run, the provider market converges to a competitive equilibrium
where provider margins approach zero economic profit (revenue equals cost
including opportunity cost of capital):

\begin{equation}
    \lim_{t \to \infty} \Pi_{\mathcal{P}, t} = 0 \quad \text{for all } \mathcal{P} \text{ in competitive equilibrium}
\end{equation}

This is the standard result for competitive markets. The Lux model achieves
this equilibrium faster than centralized markets because: (i)~provider
switching costs are near zero, removing the friction that sustains
supranormal profits; (ii)~reputation is public and verifiable, reducing
information asymmetry; and (iii)~market clearing is automated and
transparent, eliminating price discrimination.

\subsection{Network Effects Without Lock-In}

The Lux provider market exhibits network effects (more providers attract
more users, which attract more providers) without the lock-in that
typically accompanies network effects in centralized platforms. The
mechanism is the open protocol: all providers implement the same
interface, so a user who joins the network to access one provider
can seamlessly use any other. The network effect accrues to the
protocol, not to any individual provider.

\section{Related Work}\label{sec:related}

Decentralized cloud computing has been explored across several
dimensions. Filecoin \cite{benet2017filecoin} provides decentralized
storage with proof-of-spacetime verification but does not address
compute, relay, or recovery services. Akash Network \cite{akash2020}
offers a decentralized compute marketplace but lacks client-side
encryption and chain-enforced SLAs. Golem \cite{golem2016} and
iExec \cite{iexec2017} provide decentralized compute but focus on
batch processing rather than real-time service markets.

The capsule architecture for user-owned encrypted data builds on
local-first software principles \cite{kleppmann2019localfirst} and
the Solid project's user-controlled data pods \cite{bernerslee2016solid}.
The Lux model extends these with chain-enforced SLAs, provider
reputation, and compute markets.

Threshold cryptography for recovery draws on Shamir's secret sharing
\cite{shamir1979share} and recent work on threshold ECDSA
\cite{gennaro2020frost} for distributed key management. The Lux
recovery provider model applies these primitives in a market context
where recovery availability is economically incentivized.

Provider reputation without surveillance relates to privacy-preserving
reputation systems studied by Androulaki et al.\ \cite{androulaki2008reputation}
and Kerschbaum \cite{kerschbaum2009reputation}. The Lux approach uses
on-chain verifiable signals rather than encrypted feedback, avoiding
the complexity of fully homomorphic reputation aggregation while
providing a practical privacy/accuracy tradeoff.

\section{Conclusion}\label{sec:conclusion}

The centralized cloud model treats providers as data custodians who
extract rent from switching costs. The Lux provider model treats
providers as service markets competing on quality, with users retaining
ownership of their keys, their data, and their exit rights.

The key contributions of this paper are:

\begin{enumerate}
    \item \textbf{Five-role provider taxonomy.} Relay, storage, gateway,
    recovery, and compute as separable, composable service roles, none of
    which require custodial access to user data.

    \item \textbf{Surveillance-free reputation.} A provider reputation
    protocol based entirely on verifiable on-chain signals---receipt
    completion, stake tenure, slashing history, proof-of-service---that
    requires no inspection of user traffic.

    \item \textbf{Chain-enforced SLAs.} Service level agreements as smart
    contracts with automatic stake-backed penalties, replacing the legal
    fiction of centralized SLA credits.

    \item \textbf{Resource markets.} Continuous double-auction markets
    for inference, storage, bandwidth, recovery, and gateway services
    with second-price settlement and EIP-1559-style dynamic pricing.

    \item \textbf{Provable portability.} A formal proof that users can
    switch providers with zero data loss and bounded downtime, even under
    Byzantine fault assumptions.
\end{enumerate}

The decentralized cloud is not a replacement for AWS. It is a different
category of infrastructure: one where the user is sovereign, the provider
is a vendor, and exit is not a migration---it is a market transaction.

\begin{thebibliography}{99}

\bibitem{synergy2024cloud}
Synergy Research Group, ``Cloud infrastructure services market reaches \$76 billion in Q4 2024,'' \emph{Synergy Research Group Report}, 2024.

\bibitem{aws2021outage}
Amazon Web Services, ``Summary of the AWS service event in the Northern Virginia (US-EAST-1) region,'' \emph{AWS Post-Event Summary}, December 2021.

\bibitem{benet2017filecoin}
J.~Benet and N.~Greco, ``Filecoin: A decentralized storage network,'' \emph{Protocol Labs Technical Report}, 2017.

\bibitem{akash2020}
G.~Gadiraju and A.~Prem, ``Akash Network: Decentralized cloud compute marketplace,'' \emph{Akash Network Whitepaper}, 2020.

\bibitem{golem2016}
The Golem Project, ``The Golem Project: Crowdfunding whitepaper,'' \emph{Golem Factory Technical Report}, 2016.

\bibitem{iexec2017}
G.~Fedak et al., ``iExec: Blockchain-based decentralized cloud computing,'' \emph{iExec Whitepaper}, 2017.

\bibitem{kleppmann2019localfirst}
M.~Kleppmann, A.~Wiggins, P.~van~Hardenberg, and M.~McGranaghan, ``Local-first software: You own your data, in spite of the cloud,'' in \emph{Onward!}, ACM, 2019.

\bibitem{bernerslee2016solid}
T.~Berners-Lee, ``Solid: A platform for decentralized social applications based on Linked Data,'' \emph{MIT CSAIL Technical Report}, 2016.

\bibitem{shamir1979share}
A.~Shamir, ``How to share a secret,'' \emph{Communications of the ACM}, vol.~22, no.~11, pp.~612--613, 1979.

\bibitem{gennaro2020frost}
R.~Gennaro and S.~Goldfeder, ``One round threshold ECDSA with identifiable abort,'' \emph{IACR ePrint 2020/540}, 2020.

\bibitem{androulaki2008reputation}
E.~Androulaki, S.~G.~Choi, S.~M.~Bellovin, and T.~Malkin, ``Reputation systems for anonymous networks,'' in \emph{PETS}, Springer, 2008.

\bibitem{kerschbaum2009reputation}
F.~Kerschbaum, ``A verifiable, centralized, coercion-free reputation system,'' in \emph{ACM WPES}, 2009.

\bibitem{roughgarden2020eip1559}
T.~Roughgarden, ``Transaction fee mechanism design for the Ethereum blockchain: An economic analysis of EIP-1559,'' \emph{arXiv preprint arXiv:2012.00854}, 2020.

\bibitem{wood2016polkadot}
G.~Wood, ``Polkadot: Vision for a heterogeneous multi-chain framework,'' \emph{Polkadot Whitepaper}, 2016.

\end{thebibliography}

\end{document}
